Forecasting chaotic dynamics from data is a standard testbed for learned dynamical models: a useful model must track a trajectory for as long as the sensitivity to initial conditions allows, and, once the trajectory has decorrelated from the truth, its free-running output should still sample the same attractor. Echo state networks (ESNs), recurrent networks in which only a linear readout is trained, were applied to chaotic time-series prediction by Jaeger and Haas~\cite{jaeger2004harnessing}, and became a reference method for this problem after Pathak et al.\ showed that a reservoir can reproduce the Lorenz and Kuramoto--Sivashinsky attractors and their Lyapunov exponents from data alone~\cite{pathak2017using}. Vlachas et al.\ compared recurrent networks trained by backpropagation through time with reservoir computing on Lorenz-96 and Kuramoto--Sivashinsky dynamics~\cite{vlachas2020backpropagation}, and the benchmark collection of Gilpin~\cite{gilpin2021chaos,gilpin2023model} compares many forecasters across a large set of chaotic systems.

The outcome of such a comparison depends on how each model family is tuned and scored; the optimiser-step study in this paper (Sec.~\ref{sec:ablations}) is a concrete example, since the gap between the ESN and a small MLP changes with the number of optimiser steps the MLP is given. This paper is a small, controlled comparison on a CPU: one data pipeline, one validation-based tuning procedure, one definition of valid prediction time (VPT) and one climate metric are applied to three forecasters (a reimplemented ESN, a one-step MLP on delay coordinates, and a small GRU~\cite{cho2014learning}) on Lorenz-63~\cite{lorenz1963deterministic} and a 10-variable Lorenz-96 system~\cite{lorenz1998optimal}. We ask (a) how the three compare in VPT and in attractor statistics, (b) how sensitive the ESN is to spectral radius and size, and (c) how much training data each needs. We state five falsifiable hypotheses in advance (Sec.~\ref{sec:setup}) and report each verdict, including those that fail. There is no new method here, and the scale (two systems, 3--5 seeds, small models) limits what can be concluded (Sec.~\ref{sec:limitations}).
